% GENERATED FILE. Edit canonical lessons, then run npm run generate:books.
% canonical-source-hash: fnv1a64:185539b1c6fd6210
% canonical-lessons: TA-C434-pulveli, TA-C434-mey, TA-C434-toppu, TA-C434-panai, TA-C434-putar

\chapter{Meadows and Groves}
\label{ch:ta-meadows-and-groves}
\hlchaptermodality{434}

\begin{chapteropening}
\textbf{By the end of this chapter:} \emph{I can talk about meadows, groves, palmyra palms and bushes, and animals grazing.}

Complete the last lesson of chapter 434: I can talk about meadows, groves, palmyra palms and bushes, and animals grazing.
\end{chapteropening}

\section[\texorpdfstring{\ta{புல்வெளி} (pulveḷi)}{புல்வெளி (pulveḷi)}]{\ta{புல்வெளி} (pulveḷi) — a meadow; grassland}

\label{lesson:TA-C434-pulveli}



\begin{quote}
Before the new one: say the Tamil for a glacier, then the Tamil for a plain.
\end{quote}

\subsection*{You'll want to know: \ta{புல்வெளி}}
\textbf{\ta{புல்வெளி}} — \emph{pulveḷi} — \textquotedblleft{}a meadow; grassland\textquotedblright{}.

\begin{quote}
\ta{புல்வெளியில்} \ta{குழந்தைகள்} \ta{ஓடி} \ta{விளையாடினார்கள்}.
\end{quote}

\begin{quote}
\emph{pulveḷiyil kuḻantaikaḷ ōṭi viḷaiyāṭiṉārkaḷ.}
\end{quote}

\begin{quote}
The children ran and played in the meadow.
\end{quote}

\ta{புல்}, grass, + \ta{வெளி}, an open space.

\begin{cousinweb}
\textbf{\ta{புல்வெளி}} (\emph{pulveḷi}) is built from \ta{புல்}, \textquotedblleft{}grass\textquotedblright{} (DEDR 4300), and \ta{வெளி}, \textquotedblleft{}an open space\textquotedblright{} (DEDR 5498).
\end{cousinweb}

\begin{grammarlens}[title={What you've built}]
One more word for getting around.
\end{grammarlens}

\subsection*{Your turn}
\begin{itemize}
\raggedright
  \item \emph{Say it:} \emph{pulveḷi}
  \item \emph{Say it:} \emph{pulveḷi}, once more
  \item \emph{Say it:} \emph{pulveḷiyil kuḻantaikaḷ ōṭi viḷaiyāṭiṉārkaḷ.} — The children ran and played in the meadow.
  \item \emph{Recall:} say the Tamil for a glacier, then the Tamil for a plain, then say \emph{pulveḷi} again
\end{itemize}

\begin{tcolorbox}[breakable,colback=teal!4,colframe=teal!35!black,title={Before you move on}]
What does \ta{புல்வெளி} mean? (\textquotedblleft{}A meadow; grassland\textquotedblright{}.) Say it once more.
\end{tcolorbox}
\section[\texorpdfstring{\ta{மேய்} (mēy)}{மேய் (mēy)}]{\ta{மேய்} (mēy) — to graze}

\label{lesson:TA-C434-mey}



\begin{quote}
Before the new one: say the Tamil for a plain, then the Tamil for a meadow.
\end{quote}

\subsection*{You'll want to know: \ta{மேய்}}
\textbf{\ta{மேய்}} — \emph{mēy} — \textquotedblleft{}to graze\textquotedblright{}.

\begin{quote}
\ta{புல்வெளியில்} \ta{பசுக்கள்} \ta{மேய்ந்தன}.
\end{quote}

\begin{quote}
\emph{pulveḷiyil pacukkaḷ mēyntaṉa.}
\end{quote}

\begin{quote}
Cows grazed in the meadow.
\end{quote}

Past (it …): \ta{மேய்ந்தது}, \emph{mēyntatu}.

To take animals out to graze is \ta{மேய்} too: \ta{ஆடு} \ta{மேய்}, to herd goats.

\begin{cousinweb}
\textbf{\ta{மேய்}} (\emph{mēy}) is native Dravidian (DEDR 5093), \textquotedblleft{}to graze, to feed\textquotedblright{}.
\end{cousinweb}

\begin{grammarlens}[title={What you've built}]
One more word for getting around.
\end{grammarlens}

\subsection*{Your turn}
\begin{itemize}
\raggedright
  \item \emph{Say it:} \emph{mēy}
  \item \emph{Say it:} \emph{mēy}, once more
  \item \emph{Say it:} \emph{pulveḷiyil pacukkaḷ mēyntaṉa.} — Cows grazed in the meadow.
  \item \emph{Recall:} say the Tamil for a plain, then the Tamil for a meadow, then say \emph{mēy} again
\end{itemize}

\begin{tcolorbox}[breakable,colback=teal!4,colframe=teal!35!black,title={Before you move on}]
What does \ta{மேய்} mean? (\textquotedblleft{}To graze\textquotedblright{}.) Say it once more.
\end{tcolorbox}
\section[\texorpdfstring{\ta{தோப்பு} (tōppu)}{தோப்பு (tōppu)}]{\ta{தோப்பு} (tōppu) — a grove; an orchard}

\label{lesson:TA-C434-toppu}



\begin{quote}
Before the new one: say the Tamil for a meadow, then the Tamil for to graze.
\end{quote}

\subsection*{You'll want to know: \ta{தோப்பு}}
\textbf{\ta{தோப்பு}} — \emph{tōppu} — \textquotedblleft{}a grove; an orchard\textquotedblright{}.

\begin{quote}
\ta{தோப்பில்} \ta{நிறைய} \ta{மாமரங்கள்} \ta{இருக்கின்றன}.
\end{quote}

\begin{quote}
\emph{tōppil niṟaiya māmaraṅkaḷ irukkiṉṟaṉa.}
\end{quote}

\begin{quote}
There are many mango trees in the grove.
\end{quote}

\ta{தென்னந்தோப்பு} is a coconut grove.

\begin{cousinweb}
\textbf{\ta{தோப்பு}} (\emph{tōppu}) is native Dravidian (DEDR 3551), \textquotedblleft{}a clump of trees, a grove\textquotedblright{}.
\end{cousinweb}

\begin{grammarlens}[title={What you've built}]
One more word for getting around.
\end{grammarlens}

\subsection*{Your turn}
\begin{itemize}
\raggedright
  \item \emph{Say it:} \emph{tōppu}
  \item \emph{Say it:} \emph{tōppu}, once more
  \item \emph{Say it:} \emph{tōppil niṟaiya māmaraṅkaḷ irukkiṉṟaṉa.} — There are many mango trees in the grove.
  \item \emph{Recall:} say the Tamil for a meadow, then the Tamil for to graze, then say \emph{tōppu} again
\end{itemize}

\begin{tcolorbox}[breakable,colback=teal!4,colframe=teal!35!black,title={Before you move on}]
What does \ta{தோப்பு} mean? (\textquotedblleft{}A grove; an orchard\textquotedblright{}.) Say it once more.
\end{tcolorbox}
\section[\texorpdfstring{\ta{பனை} (paṉai)}{பனை (paṉai)}]{\ta{பனை} (paṉai) — a palmyra palm}

\label{lesson:TA-C434-panai}



\begin{quote}
Before the new one: say the Tamil for to graze, then the Tamil for a grove.
\end{quote}

\subsection*{You'll want to know: \ta{பனை}}
\textbf{\ta{பனை}} — \emph{paṉai} — \textquotedblleft{}a palmyra palm\textquotedblright{}.

\begin{quote}
\ta{கடற்கரையோரம்} \ta{உயரமான} \ta{பனைமரங்கள்} \ta{நிற்கின்றன}.
\end{quote}

\begin{quote}
\emph{kaṭaṟkaraiyōram uyaramāṉa paṉaimaraṅkaḷ niṟkiṉṟaṉa.}
\end{quote}

\begin{quote}
Tall palmyra palms stand along the shore.
\end{quote}

\ta{பனைமரம்} is the whole tree; \ta{பனையோலை}, its leaf.

\begin{cousinweb}
\textbf{\ta{பனை}} (\emph{paṉai}) is native Dravidian (DEDR 4037), \textquotedblleft{}the palmyra palm\textquotedblright{}.
\end{cousinweb}

\begin{grammarlens}[title={What you've built}]
One more word for getting around.
\end{grammarlens}

\subsection*{Your turn}
\begin{itemize}
\raggedright
  \item \emph{Say it:} \emph{paṉai}
  \item \emph{Say it:} \emph{paṉai}, once more
  \item \emph{Say it:} \emph{kaṭaṟkaraiyōram uyaramāṉa paṉaimaraṅkaḷ niṟkiṉṟaṉa.} — Tall palmyra palms stand along the shore.
  \item \emph{Recall:} say the Tamil for to graze, then the Tamil for a grove, then say \emph{paṉai} again
\end{itemize}

\begin{tcolorbox}[breakable,colback=teal!4,colframe=teal!35!black,title={Before you move on}]
What does \ta{பனை} mean? (\textquotedblleft{}A palmyra palm\textquotedblright{}.) Say it once more.
\end{tcolorbox}
\section[\texorpdfstring{\ta{புதர்} (putar)}{புதர் (putar)}]{\ta{புதர்} (putar) — a bush; a thicket}

\label{lesson:TA-C434-putar}



\begin{quote}
Before the new one: say the Tamil for a grove, then the Tamil for a palmyra palm.
\end{quote}

\subsection*{You'll want to know: \ta{புதர்}}
\textbf{\ta{புதர்}} — \emph{putar} — \textquotedblleft{}a bush; a thicket\textquotedblright{}.

\begin{quote}
\ta{புதரில்} \ta{ஒரு} \ta{பாம்பு} \ta{மறைந்திருந்தது}.
\end{quote}

\begin{quote}
\emph{putaril oru pāmpu maṟaintiruntatu.}
\end{quote}

\begin{quote}
A snake was hiding in the bush.
\end{quote}

The plural is \ta{புதர்கள்}.

\begin{cousinweb}
\textbf{\ta{புதர்}} (\emph{putar}) is native Dravidian (DEDR 4509), \textquotedblleft{}a bush, low jungle\textquotedblright{}.
\end{cousinweb}

\begin{grammarlens}[title={What you've built}]
One more word for getting around.
\end{grammarlens}

\subsection*{Your turn}
\begin{itemize}
\raggedright
  \item \emph{Say it:} \emph{putar}
  \item \emph{Say it:} \emph{putar}, once more
  \item \emph{Say it:} \emph{putaril oru pāmpu maṟaintiruntatu.} — A snake was hiding in the bush.
  \item \emph{Recall:} say the Tamil for a grove, then the Tamil for a palmyra palm, then say \emph{putar} again
\end{itemize}

\begin{tcolorbox}[breakable,colback=teal!4,colframe=teal!35!black,title={Before you move on}]
What does \ta{புதர்} mean? (\textquotedblleft{}A bush; a thicket\textquotedblright{}.) Say it once more.
\end{tcolorbox}
